\section{Conclusiones}

El planificador implementado traduce la política de \textit{proportional-share
scheduling} de Waldspurger y Weihl a un mecanismo ejecutable y verificable en
espacio de usuario, cumpliendo las restricciones de concurrencia del
enunciado: exclusión mutua verificada en tiempo de ejecución (el
\texttt{assert} de \texttt{sync\_dispatch}), sin espera activa (todo bloqueo
pasa por variables de condición con predicado explícito), y sin un único
candado que envuelva la ejecución del trabajo real. Los siete casos mínimos
de funcionamiento cumplen, verificados bajo AddressSanitizer,
UndefinedBehaviorSanitizer, LeakSanitizer y ThreadSanitizer sin hallazgos, y
el experimento de proporcionalidad confirma que el sorteo pondera
correctamente por boletos dentro de una ventana de observación truncada,
con error absoluto medio muy por debajo de la tolerancia exigida.

La extensión elegida, \textit{compensation tickets}, reutiliza el mecanismo
cooperativo/quantum ya construido para modelar una tarea que cede
tempranamente y recupera participación de CPU mediante boletos inflados
temporalmente, sin alterar el comportamiento del planificador base cuando la
extensión no está activa. La cuantificación estadística final de cuánta
participación recupera --- y su sensibilidad al tamaño de la ventana
observada --- se completa en la Sección~\ref{sec:resultados-extension}.

La comparación con el paper original deja claras las simplificaciones
deliberadas de este prototipo: un único CPU lógico simulado en vez de
integración real con un \textit{kernel}, una carga de trabajo sintética e
instrumentada en vez de programas reales, y un único tipo de recurso en vez
de la generalización a memoria y red que demuestra Mach 3.0. Estas
simplificaciones son coherentes con el objetivo del proyecto --- verificar
el protocolo de sincronización y la política de asignación, no medir
rendimiento de un sistema real --- y se documentan explícitamente en vez de
presentarse como limitaciones no reconocidas.
